% Options for packages loaded elsewhere
\PassOptionsToPackage{unicode}{hyperref}
\PassOptionsToPackage{hyphens}{url}
\documentclass[
  brazilian,
  11pt,
]{article}
\usepackage{xcolor}
\usepackage[margin=0.88in]{geometry}
\usepackage{amsmath,amssymb}
\setcounter{secnumdepth}{5}
\usepackage{iftex}
\ifPDFTeX
  \usepackage[T1]{fontenc}
  \usepackage[utf8]{inputenc}
  \usepackage{textcomp} % provide euro and other symbols
\else % if luatex or xetex
  \usepackage{unicode-math} % this also loads fontspec
  \defaultfontfeatures{Scale=MatchLowercase}
  \defaultfontfeatures[\rmfamily]{Ligatures=TeX,Scale=1}
\fi
\usepackage{lmodern}
\ifPDFTeX\else
  % xetex/luatex font selection
\fi
% Use upquote if available, for straight quotes in verbatim environments
\IfFileExists{upquote.sty}{\usepackage{upquote}}{}
\IfFileExists{microtype.sty}{% use microtype if available
  \usepackage[]{microtype}
  \UseMicrotypeSet[protrusion]{basicmath} % disable protrusion for tt fonts
}{}
\makeatletter
\@ifundefined{KOMAClassName}{% if non-KOMA class
  \IfFileExists{parskip.sty}{%
    \usepackage{parskip}
  }{% else
    \setlength{\parindent}{0pt}
    \setlength{\parskip}{6pt plus 2pt minus 1pt}}
}{% if KOMA class
  \KOMAoptions{parskip=half}}
\makeatother
\usepackage{graphicx}
\makeatletter
\newsavebox\pandoc@box
\newcommand*\pandocbounded[1]{% scales image to fit in text height/width
  \sbox\pandoc@box{#1}%
  \Gscale@div\@tempa{\textheight}{\dimexpr\ht\pandoc@box+\dp\pandoc@box\relax}%
  \Gscale@div\@tempb{\linewidth}{\wd\pandoc@box}%
  \ifdim\@tempb\p@<\@tempa\p@\let\@tempa\@tempb\fi% select the smaller of both
  \ifdim\@tempa\p@<\p@\scalebox{\@tempa}{\usebox\pandoc@box}%
  \else\usebox{\pandoc@box}%
  \fi%
}
% Set default figure placement to htbp
\def\fps@figure{htbp}
\makeatother
% definitions for citeproc citations
\NewDocumentCommand\citeproctext{}{}
\NewDocumentCommand\citeproc{mm}{%
  \begingroup\def\citeproctext{#2}\cite{#1}\endgroup}
\makeatletter
 % allow citations to break across lines
 \let\@cite@ofmt\@firstofone
 % avoid brackets around text for \cite:
 \def\@biblabel#1{}
 \def\@cite#1#2{{#1\if@tempswa , #2\fi}}
\makeatother
\newlength{\cslhangindent}
\setlength{\cslhangindent}{1.5em}
\newlength{\csllabelwidth}
\setlength{\csllabelwidth}{3em}
\newenvironment{CSLReferences}[2] % #1 hanging-indent, #2 entry-spacing
 {\begin{list}{}{%
  \setlength{\itemindent}{0pt}
  \setlength{\leftmargin}{0pt}
  \setlength{\parsep}{0pt}
  % turn on hanging indent if param 1 is 1
  \ifodd #1
   \setlength{\leftmargin}{\cslhangindent}
   \setlength{\itemindent}{-1\cslhangindent}
  \fi
  % set entry spacing
  \setlength{\itemsep}{#2\baselineskip}}}
 {\end{list}}
\usepackage{calc}
\newcommand{\CSLBlock}[1]{\hfill\break\parbox[t]{\linewidth}{\strut\ignorespaces#1\strut}}
\newcommand{\CSLLeftMargin}[1]{\parbox[t]{\csllabelwidth}{\strut#1\strut}}
\newcommand{\CSLRightInline}[1]{\parbox[t]{\linewidth - \csllabelwidth}{\strut#1\strut}}
\newcommand{\CSLIndent}[1]{\hspace{\cslhangindent}#1}
\ifLuaTeX
\usepackage[bidi=basic,shorthands=off,]{babel}
\else
\usepackage[bidi=default,shorthands=off,]{babel}
\fi
\ifLuaTeX
  \usepackage{selnolig} % disable illegal ligatures
\fi
\setlength{\emergencystretch}{3em} % prevent overfull lines
\providecommand{\tightlist}{%
  \setlength{\itemsep}{0pt}\setlength{\parskip}{0pt}}
\usepackage{booktabs}
\usepackage{longtable}
\usepackage{array}
\usepackage{float}
\usepackage{xcolor}
\usepackage{hyperref}
\usepackage{caption}
\definecolor{deepblue}{HTML}{17365D}
\definecolor{softblue}{HTML}{EAF1F8}
\definecolor{darkgray}{HTML}{404040}
\hypersetup{colorlinks=true,linkcolor=deepblue,citecolor=deepblue,urlcolor=deepblue}
\Urlmuskip=0mu plus 1mu\relax
\setlength{\parindent}{0pt}
\setlength{\parskip}{0.48em}
\widowpenalty=10000
\clubpenalty=10000
\captionsetup{font=small,labelfont=bf}
\usepackage{booktabs}
\usepackage{longtable}
\usepackage{array}
\usepackage{multirow}
\usepackage{wrapfig}
\usepackage{float}
\usepackage{colortbl}
\usepackage{pdflscape}
\usepackage{tabu}
\usepackage{threeparttable}
\usepackage{threeparttablex}
\usepackage[normalem]{ulem}
\usepackage{makecell}
\usepackage{xcolor}
\usepackage{bookmark}
\IfFileExists{xurl.sty}{\usepackage{xurl}}{} % add URL line breaks if available
\urlstyle{same}
\hypersetup{
  pdftitle={Public cues, SDiD individuais e modelos de fatores latentes},
  pdfauthor={Nota diagnóstica},
  pdflang={pt-BR},
  hidelinks,
  pdfcreator={LaTeX via pandoc}}

\title{Public cues, SDiD individuais e modelos de fatores latentes}
\usepackage{etoolbox}
\makeatletter
\providecommand{\subtitle}[1]{% add subtitle to \maketitle
  \apptocmd{\@title}{\par {\large #1 \par}}{}{}
}
\makeatother
\subtitle{Chile, Uruguai, Gabão, Catar, Austrália e a sensibilidade com
Brasil}
\author{Nota diagnóstica}
\date{22 de setembro de 2026}

\begin{document}
\maketitle

{
\setcounter{tocdepth}{1}
\tableofcontents
}
\begin{center}
\small\color{darkgray}
\textbf{Status:} análise diagnóstica fora do pipeline \texttt{targets}. O relatório importa resultados já estimados; sua compilação não reestima modelos.
\end{center}

\begin{abstract}
Este relatório apresenta seis estimativas individuais de synthetic difference-in-differences (SDiD) e quatro modelos pooled de fatores latentes para avaliar mudanças na distância dos votos na Assembleia Geral das Nações Unidas em relação à China após cues públicos de posição comercial. Os SDiD individuais para Chile, Uruguai, Gabão, Catar e Austrália são imprecisos e todos os intervalos de 95\% incluem zero. Nos modelos pooled com os cinco casos, as estimativas também ficam próximas de zero. Quando Gabão e Catar são excluídos e o Brasil é incluído, o ATT torna-se mais negativo. A especificação em que controles entram e saem do conjunto de comparação produz ATT de -0,1359, erro-padrão bootstrap de 0,0657 e intervalo de 95\% [-0,2646; -0,0073]. Esse resultado é sugestivo de aproximação à China, mas a inferência permanece frágil porque há somente quatro países tratados.
\end{abstract}

\section{Resultado principal}\label{resultado-principal}

O padrão empírico é simples. As estimativas individuais e os modelos
pooled com os cinco casos não oferecem evidência precisa de mudança. A
sensibilidade que mantém Austrália, Brasil, Chile e Uruguai como casos
tratados produz estimativas mais negativas:

\begin{itemize}
\tightlist
\item
  com controles limpos, ATT = -0,1158, IC 95\% {[}-0,2441; 0,0124{]},
  \(p=\) 0,077;
\item
  com controles entrando e saindo do conjunto de comparação, ATT =
  -0,1359, IC 95\% {[}-0,2646; -0,0073{]}, \(p=\) 0,038.
\end{itemize}

Como o outcome é a distância absoluta à China, valores negativos indicam
aproximação. A diferença entre as especificações de cinco e quatro casos
mostra que Gabão e Catar alteram substantivamente o resultado pooled.
Ela não identifica, por si só, qual caso gera a mudança nem transforma a
análise em evidência causal definitiva.

\section{Desenho da análise}\label{desenho-da-anuxe1lise}

\subsection{Unidade, período e
outcome}\label{unidade-peruxedodo-e-outcome}

A unidade é o país-ano. A janela comum vai de 1997 a 2022. O outcome é a
distância absoluta entre o ponto ideal do país e o ponto ideal da China
nas votações da Assembleia Geral das Nações Unidas. A análise termina em
2022 porque esse é o último ano com a medida de ranking comercial usada
para determinar a elegibilidade dos controles.

O input empírico é o target armazenado
\texttt{china\_top\_m2\_goods\_panel}, hash \texttt{7907d353f7750752},
produzido em 2026-05-24 16:31:38.188157. O relatório e os scripts
diagnósticos apenas leram esse objeto. Nenhum comando
\texttt{targets::tar\_make()} foi executado.

\subsection{Casos tratados e timing}\label{casos-tratados-e-timing}

O conjunto original contém cinco casos com cue público codificado:

\begin{itemize}
\tightlist
\item
  Chile: 2008;
\item
  Uruguai: 2013;
\item
  Gabão: 2017;
\item
  Catar: 2021;
\item
  Austrália: 2009.
\end{itemize}

Para a Austrália, o SDiD também é estimado com início em 2010, conforme
a alternativa substantiva solicitada. A sensibilidade pooled sem Gabão e
Catar inclui o Brasil, tratado a partir de 2009. Assim, o grupo reduzido
é Austrália, Brasil, Chile e Uruguai.

Nos modelos de fatores latentes, o tratamento é

\[
D_{it}
=
\mathbf{1}\{i\text{ é caso focal}\}
\mathbf{1}\{t\geq \text{cue}_i\}
\mathbf{1}\{\text{China é top 1}_{it}\}.
\]

Países não focais nunca recebem \(D_{it}=1\).

\subsection{Duas regras para o conjunto de
comparação}\label{duas-regras-para-o-conjunto-de-comparauxe7uxe3o}

Cada grupo focal é estimado sob duas regras:

\begin{enumerate}
\def\labelenumi{\arabic{enumi}.}
\tightlist
\item
  \textbf{Controles limpos.} O donor pool contém 105 países com painel
  completo que nunca têm a China como maior destino de exportações em
  1997--2022.
\item
  \textbf{Controles que entram e saem.} O modelo começa com os 159
  painéis completos. Para países não focais, anos China-top são
  mascarados e não entram como observações de controle. Quando a China
  deixa de ser top 1, o país volta ao conjunto de comparação. Esses
  países não são tratados adicionais.
\end{enumerate}

A Mongólia ilustra a diferença. Ela recebe \(D_{it}=0\) em todos os
anos, mas tem a China como top 1 durante toda a janela. Seus 26
país-anos são, portanto, inelegíveis como controle. Como não resta
observação utilizável, o \texttt{fect} a remove do ajuste; ela não entra
no ATT focal.

\section{Estimadores e inferência}\label{estimadores-e-inferuxeancia}

\subsection{SDiD individuais}\label{sdid-individuais}

Cada SDiD contém um país tratado e os mesmos 105 controles limpos. O
estimador combina ponderação de unidades e de períodos para construir o
contrafactual (\citeproc{ref-arkhangelsky_etal2021}{Arkhangelsky et al.
2021}). Os erros-padrão placebo usam 5000 reestimações por ajuste e seed
20260520. Os \(p\)-valores apresentados são aproximações normais
baseadas nesses erros-padrão, não \(p\)-valores exatos de randomização.

\subsection{Fatores latentes pooled}\label{fatores-latentes-pooled}

Os modelos pooled usam o estimador de efeitos fixos interativos do
pacote \texttt{fect 2.4.5}, com efeitos fixos bidirecionais e seleção do
número de fatores por validação cruzada entre \(r=0,1,2,3\)
(\citeproc{ref-xu2017gsynth}{Xu 2017};
\citeproc{ref-liu2024practical}{Liu, Wang, e Xu 2024}). Cada ajuste usa
1000 repetições de bootstrap por unidade, seed 42 e \texttt{min.T0 = 1}.
Nos modelos com controles que entram e saem, células inelegíveis são
tratadas como outcomes não observados e estimadas pelo procedimento de
painel incompleto do \texttt{fect}.

\section{Resultados dos fatores
latentes}\label{resultados-dos-fatores-latentes}

\begin{table}[!h]
\centering
\caption{\label{tab:table-ife-main}Modelos pooled de fatores latentes para a distância ideal à China.}
\centering
\resizebox{\ifdim\width>\linewidth\linewidth\else\width\fi}{!}{
\fontsize{8.5}{10.5}\selectfont
\begin{tabular}[t]{>{\raggedright\arraybackslash}p{6.0cm}rrcrrc}
\toprule
Especificação & ATT & EP bootstrap & IC 95\% & p-valor & Fatores (r) & N tratados/controles\\
\midrule
Cinco casos + controles limpos & -0,0236 & 0,0750 & {}[-0,1706; 0,1234] & 0,753 & 1 & 5/105\\
Cinco casos + controles que entram/saem & -0,0213 & 0,0618 & {}[-0,1425; 0,1000] & 0,731 & 1 & 5/153\\
Austrália, Brasil, Chile e Uruguai + controles limpos & -0,1158 & 0,0654 & {}[-0,2441; 0,0124] & 0,077 & 0 & 4/105\\
Austrália, Brasil, Chile e Uruguai + controles que entram/saem & -0,1359 & 0,0657 & {}[-0,2646; -0,0073] & 0,038 & 0 & 4/152\\
\bottomrule
\end{tabular}}
\end{table}

\textit{Nota:} ATT negativo indica redução da distância à China. Os
intervalos usam a aproximação normal com erro-padrão obtido por
bootstrap de unidades. O número de fatores é selecionado pela regra de
um erro-padrão da validação cruzada.

\begin{figure}

{\centering \includegraphics[width=0.96\linewidth]{/Users/manoelgaldino/Documents/DCP/Papers/RDD Trade/red_trade/quality_reports/selected_public_cue_sdid_fect/report_assets/figure_1_ife_estimates} 

}

\caption{Estimativas pooled de fatores latentes. As barras mostram intervalos de 95\%. A linha tracejada marca efeito zero. Em todos os modelos, somente os casos focais recebem tratamento; na alternativa ampliada, países não focais entram e saem apenas do conjunto de comparação.}\label{fig:figure-ife}
\end{figure}

Os dois modelos com cinco casos produzem praticamente o mesmo ATT:
-0,0236 com controles limpos e -0,0213 com controles que entram e saem.
Ambos selecionam um fator latente e têm intervalos amplos em torno de
zero.

Nos modelos sem Gabão e Catar, com inclusão do Brasil, a validação
cruzada seleciona \(r=0\). A estimativa com controles limpos é negativa,
mas seu intervalo ainda alcança zero. Ao permitir que controles
contaminados saiam apenas durante seus anos China-top, a estimativa é
-0,1359 e o limite superior do intervalo é -0,0073.

\begin{table}[!h]
\centering
\caption{\label{tab:table-ife-composition}Composição dos grupos focais e alternância do conjunto de comparação.}
\centering
\resizebox{\ifdim\width>\linewidth\linewidth\else\width\fi}{!}{
\fontsize{8.5}{10.5}\selectfont
\begin{tabular}[t]{>{\raggedright\arraybackslash}p{5.2cm}lrr}
\toprule
Especificação & Casos focais & Controles que alternam & País-anos mascarados\\
\midrule
Cinco casos + controles limpos & AUS, CHL, GAB, QAT, URY & 0 & 0\\
Cinco casos + controles que entram/saem & AUS, CHL, GAB, QAT, URY & 48 & 448\\
Austrália, Brasil, Chile e Uruguai + controles limpos & AUS, BRA, CHL, URY & 0 & 0\\
Austrália, Brasil, Chile e Uruguai + controles que entram/saem & AUS, BRA, CHL, URY & 47 & 434\\
\bottomrule
\end{tabular}}
\end{table}

No modelo ampliado de cinco casos, 48 controles alternam sua
elegibilidade e 448 país-anos são mascarados. Na versão com quatro
casos, são 47 controles e 434 país-anos. Essas células não geram ATTs
adicionais.

\section{Resultados dos SDiD
individuais}\label{resultados-dos-sdid-individuais}

\begin{table}[H]
\centering
\caption{\label{tab:table-sdid}SDiD individuais com erros-padrão placebo.}
\centering
\resizebox{\ifdim\width>\linewidth\linewidth\else\width\fi}{!}{
\fontsize{8.2}{10.2}\selectfont
\begin{tabular}[t]{llrrcrcr}
\toprule
País & Timing & ATT & EP placebo & IC 95\% & p-valor & Pré/Pós & Doadores\\
\midrule
Chile & Cue público em 2008 & -0,1301 & 0,1545 & {}[-0,4328; 0,1726] & 0,400 & 11/15 & 105\\
Uruguai & Cue público em 2013 & -0,0682 & 0,1514 & {}[-0,3649; 0,2285] & 0,652 & 16/10 & 105\\
Gabão & Cue público em 2017 & -0,0429 & 0,1175 & {}[-0,2732; 0,1874] & 0,715 & 20/6 & 105\\
Catar & Cue público em 2021 & 0,1263 & 0,1221 & {}[-0,1131; 0,3657] & 0,301 & 24/2 & 105\\
Austrália & Cue público em 2009 & 0,0420 & 0,1545 & {}[-0,2608; 0,3449] & 0,786 & 12/14 & 105\\
\addlinespace
Austrália & Timing alternativo em 2010 & 0,0900 & 0,1468 & {}[-0,1978; 0,3778] & 0,540 & 13/13 & 105\\
\bottomrule
\end{tabular}}
\end{table}

\begin{figure}[H]

{\centering \includegraphics[width=0.96\linewidth]{/Users/manoelgaldino/Documents/DCP/Papers/RDD Trade/red_trade/quality_reports/selected_public_cue_sdid_fect/report_assets/figure_2_sdid_estimates} 

}

\caption{Estimativas SDiD individuais. As barras mostram intervalos de 95\% construídos com erros-padrão placebo de 5.000 reestimações. O ajuste australiano de 2010 é uma alternativa de timing; os demais marcam o cue público selecionado.}\label{fig:figure-sdid}
\end{figure}

Chile, Uruguai e Gabão têm estimativas negativas; Catar e os dois
timings australianos têm estimativas positivas. Nenhum intervalo
individual exclui zero. O resultado do Catar deve ser interpretado com
cautela adicional porque dispõe de apenas dois anos pós-tratamento.

\begin{table}[H]
\centering
\caption{\label{tab:table-sdid-fit}Diagnósticos de ajuste pré e pós-tratamento dos SDiD.}
\centering
\resizebox{\ifdim\width>\linewidth\linewidth\else\width\fi}{!}{
\fontsize{8.5}{10.5}\selectfont
\begin{tabular}[t]{lrrrrr}
\toprule
Ajuste & RMSPE pré & RMSPE pós & Razão pós/pré & Anos pré & Anos pós\\
\midrule
Chile, 2008 & 0,0799 & 0,1843 & 2,306 & 11 & 15\\
Uruguai, 2013 & 0,1161 & 0,1705 & 1,468 & 16 & 10\\
Gabão, 2017 & 0,1108 & 0,1268 & 1,145 & 20 & 6\\
Catar, 2021 & 0,1311 & 0,1675 & 1,278 & 24 & 2\\
Austrália, 2009 & 0,1546 & 0,1210 & 0,783 & 12 & 14\\
\addlinespace
Austrália, 2010 & 0,1511 & 0,1474 & 0,976 & 13 & 13\\
\bottomrule
\end{tabular}}
\end{table}

A razão RMSPE pós/pré é maior que um para Chile, Uruguai, Gabão e Catar,
mas a precisão placebo permanece baixa. Na Austrália, a razão é inferior
a um em 2009 e próxima de um em 2010, consistente com a ausência de uma
separação pós-tratamento clara.

\clearpage

\section{Austrália: sensibilidade do SDiD a 2006, 2007 e
2008}\label{austruxe1lia-sensibilidade-do-sdid-a-2006-2007-e-2008}

Esta seção acrescenta três ajustes pontuais já estimados, com início de
tratamento imposto em 2006, 2007 e 2008. Todos mantêm a janela
1997--2022 e os mesmos 105 controles limpos. Essas datas são
alternativas de sensibilidade e não substituem a codificação do cue
público em 2009. Não foram calculados erros-padrão, intervalos de
confiança ou testes placebo para esses três ajustes.

As figuras reproduzem o método nativo
\texttt{plot.synthdid\_estimate()}, chamado com
\texttt{se.method = "none"}. O ATT compara a mudança da diferença entre
Austrália e controle sintético, usando os pesos temporais do estimador;
a diferença bruta de níveis entre as linhas não é o ATT. Alterar o
início do tratamento modifica a divisão pré/pós e reestima os pesos de
unidades e períodos.

\subsection{Tratamento imposto em
2006}\label{tratamento-imposto-em-2006}

ATT = 0,2077, com nove anos pré-tratamento (1997--2005) e 17 anos
tratados (2006--2022).

\begin{figure}[H]

{\centering \includegraphics[width=1\linewidth]{/Users/manoelgaldino/Documents/DCP/Papers/RDD Trade/red_trade/quality_reports/selected_public_cue_sdid_fect/figures/australia_sdid_2006_native} 

}

\caption{Austrália: gráfico nativo do SDiD com início imposto em 2006. Outcome: distância absoluta do ponto ideal à China. Janela 1997--2022, 105 doadores limpos. A faixa inferior apresenta os pesos temporais; a linha vertical marca o último ano pré-tratamento. Estimativa pontual, sem inferência.}\label{fig:figure-australia-2006}
\end{figure}

\clearpage

\subsection{Tratamento imposto em
2007}\label{tratamento-imposto-em-2007}

ATT = 0,1285, com dez anos pré-tratamento (1997--2006) e 16 anos
tratados (2007--2022).

\begin{figure}[H]

{\centering \includegraphics[width=1\linewidth]{/Users/manoelgaldino/Documents/DCP/Papers/RDD Trade/red_trade/quality_reports/selected_public_cue_sdid_fect/figures/australia_sdid_2007_native} 

}

\caption{Austrália: gráfico nativo do SDiD com início imposto em 2007. Outcome: distância absoluta do ponto ideal à China. Janela 1997--2022, 105 doadores limpos. A faixa inferior apresenta os pesos temporais; a linha vertical marca o último ano pré-tratamento. Estimativa pontual, sem inferência.}\label{fig:figure-australia-2007}
\end{figure}

\clearpage

\subsection{Tratamento imposto em
2008}\label{tratamento-imposto-em-2008}

ATT = 0,0701, com 11 anos pré-tratamento (1997--2007) e 15 anos tratados
(2008--2022).

\begin{figure}[H]

{\centering \includegraphics[width=1\linewidth]{/Users/manoelgaldino/Documents/DCP/Papers/RDD Trade/red_trade/quality_reports/selected_public_cue_sdid_fect/figures/australia_sdid_2008_native} 

}

\caption{Austrália: gráfico nativo do SDiD com início imposto em 2008. Outcome: distância absoluta do ponto ideal à China. Janela 1997--2022, 105 doadores limpos. A faixa inferior apresenta os pesos temporais; a linha vertical marca o último ano pré-tratamento. Estimativa pontual, sem inferência.}\label{fig:figure-australia-2008}
\end{figure}

Os três ATTs são positivos, indicando maior distância à China em relação
ao contrafactual estimado. Sem inferência para esses ajustes, os sinais
não demonstram um efeito diferente de zero. A variação entre datas
também não constitui três testes independentes.

\clearpage

\section{Auditoria do timing e do
tratamento}\label{auditoria-do-timing-e-do-tratamento}

\begin{table}[!h]
\centering
\caption{\label{tab:table-timing}Timing e duração do tratamento nos modelos pooled.}
\centering
\resizebox{\ifdim\width>\linewidth\linewidth\else\width\fi}{!}{
\fontsize{7.3}{9.3}\selectfont
\begin{tabular}[t]{llrrrrrrr}
\toprule
Escopo & País & Ano do cue & Primeiro ano tratado & Último ano tratado & Períodos tratados & Entradas & Saídas & Anos China-top antes do cue\\
\midrule
Cinco casos & Chile & 2008 & 2008 & 2022 & 15 & 1 & 0 & 1\\
Cinco casos & Austrália & 2009 & 2009 & 2022 & 14 & 1 & 0 & 0\\
Cinco casos & Uruguai & 2013 & 2013 & 2022 & 10 & 1 & 0 & 0\\
Cinco casos & Gabão & 2017 & 2017 & 2022 & 6 & 1 & 0 & 0\\
Cinco casos & Catar & 2021 & 2021 & 2022 & 2 & 1 & 0 & 0\\
\addlinespace
Quatro casos sem Gabão/Catar, com Brasil & Chile & 2008 & 2008 & 2022 & 15 & 1 & 0 & 1\\
Quatro casos sem Gabão/Catar, com Brasil & Austrália & 2009 & 2009 & 2022 & 14 & 1 & 0 & 0\\
Quatro casos sem Gabão/Catar, com Brasil & Brasil & 2009 & 2009 & 2022 & 14 & 1 & 0 & 0\\
Quatro casos sem Gabão/Catar, com Brasil & Uruguai & 2013 & 2013 & 2022 & 10 & 1 & 0 & 0\\
\bottomrule
\end{tabular}}
\end{table}

Todos os casos focais têm uma entrada e nenhuma saída entre 1997 e 2022.
O Chile possui um ano China-top anterior ao cue codificado; esse ano
permanece não tratado, pois o tratamento exige simultaneamente status
China-top e cue público. A ausência de saídas entre os focais significa
que a alternância no segundo donor pool ocorre somente entre controles.

\section{Interpretação e limites}\label{interpretauxe7uxe3o-e-limites}

Quatro conclusões são sustentadas pelos resultados:

\begin{enumerate}
\def\labelenumi{\arabic{enumi}.}
\tightlist
\item
  \textbf{Os SDiD individuais são imprecisos.} Nenhum dos seis
  intervalos exclui zero, inclusive os dois timings australianos.
\item
  \textbf{O pooled de cinco casos é próximo de zero.} A ampliação do
  conjunto de comparação reduz o erro-padrão, mas não muda a conclusão
  substantiva.
\item
  \textbf{A exclusão de Gabão e Catar é substantiva.} Com Brasil,
  Austrália, Chile e Uruguai, os dois ATTs tornam-se mais negativos. O
  modelo com controles entrando e saindo tem intervalo que não inclui
  zero sob a aproximação bootstrap utilizada.
\item
  \textbf{A inferência dos quatro casos é frágil.} O \texttt{fect} não
  calculou o teste F por insuficiência de unidades tratadas. Um
  \(p\)-valor abaixo de 0,05 em uma especificação não elimina a
  sensibilidade a composição, timing e poucos clusters tratados.
\end{enumerate}

O relatório também tem limites de desenho. A codificação de cues
públicos é selecionada e não constitui um levantamento multinacional
exaustivo. Os modelos são outcome-only e não incluem covariáveis
adicionais. O uso de outcomes mascarados para controlar contaminação
depende da capacidade do modelo de fatores latentes de recuperar a
estrutura comum a partir das células observadas. Por fim, a Austrália
exige transparência sobre a distinção entre o cue em 2009 e o timing
alternativo de 2010.

\section{Reprodutibilidade e scripts
usados}\label{reprodutibilidade-e-scripts-usados}

\begin{table}[!h]
\centering
\caption{\label{tab:table-scripts}Scripts e fontes usados na produção da análise e do relatório.}
\centering
\fontsize{8.5}{10.5}\selectfont
\begin{tabular}[t]{>{\raggedright\arraybackslash}p{3.2cm}>{\raggedright\arraybackslash}p{10.7cm}}
\toprule
Componente & Função\\
\midrule
Estimação & Estima os seis SDiD e os quatro modelos IFE; grava auditorias e manifests.\\
Placebos SDiD & Implementa os placebos SDiD determinísticos e checkpointados.\\
Assets do relatório & Prepara as tabelas, figuras e metadados deste relatório; não reestima modelos.\\
Fonte do relatório & Importa os assets e apresenta método, resultados, limitações e proveniência.\\
Sensibilidade australiana & Estima o SDiD pontual australiano para o ano indicado, sem inferência.\\
\addlinespace
Gráficos nativos australianos & Exporta o gráfico nativo de cada ajuste australiano salvo, sem inferência.\\
\bottomrule
\end{tabular}
\end{table}

Os arquivos usados são:

\begin{itemize}
\item \nolinkurl{scripts/diagnostics/estimate_selected_public_cue_sdid_fect.R}
\item \nolinkurl{scripts/diagnostics/sdid_placebo_helpers.R}
\item \nolinkurl{scripts/diagnostics/build_selected_public_cue_report_assets.R}
\item \nolinkurl{reports/selected_public_cue_sdid_fect/relatorio_sdid_fatores_latentes.Rmd}
\item \nolinkurl{scripts/diagnostics/estimate_australia_sdid_2007_point.R}
\item \nolinkurl{scripts/diagnostics/plot_australia_sdid_2006_native.R}
\end{itemize}

Os dois scripts australianos usam
\texttt{AUSTRALIA\_SDID\_TREATMENT\_YEAR}, apesar dos anos em seus
nomes. Para reproduzir cada alternativa, defina essa variável como 2006,
2007 ou 2008 e execute primeiro a estimação pontual e depois a
exportação do gráfico. Esta atualização apenas importa os resultados
existentes.

O script de estimação lê o target armazenado
\texttt{china\_top\_m2\_goods\_panel}; ele não executa o pipeline. O
helper de placebos implementa as 5000 reestimações checkpointadas de
cada SDiD. O script de assets lê somente CSVs e manifests concluídos e
produz as tabelas e figuras deste documento. Os hashes SHA-256 completos
constam do CSV de proveniência:

\nolinkurl{quality_reports/selected_public_cue_sdid_fect/report_assets/table_5_scripts.csv}

As versões centrais registradas na execução são \texttt{fect 2.4.5} e
\texttt{synthdid 0.0.9}. Os quatro modelos IFE usam 1000 bootstraps e os
seis SDiD usam, ao todo, 30000 reestimações placebo verificadas.

\section{Arquivos produzidos}\label{arquivos-produzidos}

Os outputs centrais são:

\begin{itemize}
\tightlist
\item
  \nolinkurl{data/processed/diagnostics/selected_public_cue_sdid_fect/fect_results.csv};
\item
  \nolinkurl{data/processed/diagnostics/selected_public_cue_sdid_fect/sdid_results.csv};
\item
  \nolinkurl{data/processed/diagnostics/selected_public_cue_sdid_fect/fect_treatment_audit.csv};
\item
  \nolinkurl{data/processed/diagnostics/selected_public_cue_sdid_fect/output_manifest.csv};
\item
  \nolinkurl{quality_reports/selected_public_cue_sdid_fect/report_assets/report_asset_manifest.csv}.
\end{itemize}

Os CSVs australianos seguem o padrão
\texttt{australia\_sdid\_ANO\_point.csv}, para 2006, 2007 e 2008, no
diretório dos resultados diagnósticos. Os gráficos vetoriais estão em:

\nolinkurl{quality_reports/selected_public_cue_sdid_fect/figures/}

Seus nomes seguem o padrão \texttt{australia\_sdid\_ANO\_native.pdf}. Os
hashes dos três CSVs e dos três gráficos foram incluídos em
\texttt{report\_source\_manifest.csv}, junto às fontes dos resultados
principais.

O relatório documenta a execução concluída em 22 de setembro de 2026.
Ele não altera o manuscrito, os targets ou os dados brutos.

\section{Referências}\label{referuxeancias}

\footnotesize

\protect\phantomsection\label{refs}
\begin{CSLReferences}{1}{0}
\bibitem[\citeproctext]{ref-arkhangelsky_etal2021}
Arkhangelsky, Dmitry, Susan Athey, David A. Hirshberg, Guido W. Imbens,
e Stefan Wager. 2021. {``Synthetic {Difference-in-Differences}''}.
\emph{American Economic Review} 111 (12): 4088--118.
\url{https://doi.org/10.1257/aer.20190159}.

\bibitem[\citeproctext]{ref-liu2024practical}
Liu, Licheng, Ye Wang, e Yiqing Xu. 2024. {``A Practical Guide to
Counterfactual Estimators for Causal Inference with Time-Series
Cross-Sectional Data''}. \emph{American Journal of Political Science} 68
(1): 160--76.

\bibitem[\citeproctext]{ref-xu2017gsynth}
Xu, Yiqing. 2017. {``Generalized Synthetic Control Method: Causal
Inference with Interactive Fixed Effects Models''}. \emph{Political
Analysis} 25 (1): 57--76. \url{https://doi.org/10.1017/pan.2016.2}.

\end{CSLReferences}

\end{document}
